% ESTILO: Vancouver (biblatex-vancouver)
% ============================================================
% cita-vancouver-config.tex
% ------------------------------------------------------------
% DESDE C1 (14/11/2026, DECISIÓN DE ALBERTO) VANCOUVER VA POR
% biblatex, COMO LOS OTROS CUATRO, CON EL ESTILO biblatex-vancouver
% DE TeX Live (vancouver.bbx v0.1, 2020/05/20). ANTES ESTE ARCHIVO
% ESTABA VACÍO A PROPÓSITO: LA IDEA ERA bibtex, QUE NO TIENE
% \parencite NI refsection, Y QUE NO SE IMPLEMENTÓ.
%   - REVISTAS: SE USA CON EL ARTÍCULO Y SU «CÓMO CITAR» (\fullcite
%     EN SU refsection), PROBADO EN C1.
%   - LIBROS: docbook-to-latex.xsl SIGUE FRENANDO CON vancouver
%     HASTA QUE SE PRUEBE ALLÁ.
% ------------------------------------------------------------
% LOS AJUSTES DE ABAJO SIGUEN EL PATRÓN DE LA GUÍA DE LA UAB
% («Citas bibliográficas según el modelo Vancouver», julio de
% 2021, en castellano): «Autor A, Autor B. Título. Revista. Año;
% volumen(número):páginas.», TODOS LOS AUTORES SEPARADOS POR COMAS
% Y «Disponible en:» ANTES DE LA DIRECCIÓN.
% ============================================================
\usepackage[
  style=vancouver,
  sortcites=true,
  backend=biber,
  backref=false,
  indexing=true,
  urldate=long,
  dateabbrev=false
]{biblatex}

% ------------------------------------------------------------
% AUTORES: TODOS CON COMA, TAMBIÉN EL ÚLTIMO («Pérez A, Gómez L»).
% babel-spanish PONE «y» ANTES DEL ÚLTIMO (MEDIDO); POR ESO VA EN
% LOS EXTRAS DEL IDIOMA, QUE SE APLICAN DESPUÉS
% ------------------------------------------------------------
% spanish.lbx NO USA finalnamedelim MIENTRAS EL CONTADOR smartand SEA
% MAYOR QUE 0 (ELIGE «y» O «e» SEGÚN EL NOMBRE): EN 0 VUELVE A
% \printdelim{finalnamedelim}, QUE ACÁ ES UNA COMA (LEÍDO EN
% spanish.lbx Y biblatex.def, TeX Live 2023)
\DeclareDelimFormat{finalnamedelim}{\addcomma\space}
\setcounter{smartand}{0}

% ------------------------------------------------------------
% NÚMERO DE LA REVISTA: vancouver.bbx IMPRIME EL CAMPO issue, PERO
% EL NÚMERO DE UN ARTÍCULO VA EN number (biblatex, Y ASÍ LO
% EXPORTA gbpublisher): SIN ESTO SALÍA «2026; 3:1-12» (MEDIDO).
% CON ESTO, «2026; 3(2):1-12»
% ------------------------------------------------------------
\DeclareFieldFormat*[article]{number}{\mkbibparens{#1}}
\renewbibmacro*{volume+number+eid}{%
  \printfield{volume}%
  \printfield{number}%
  \printfield{issue}%
  \printfield{eid}}

% ------------------------------------------------------------
% PUNTO FINAL: vancouver.bbx LO QUITA (\finentrypunct VACÍO). LA GUÍA
% CIERRA CADA REFERENCIA CON PUNTO, SALVO LAS QUE TERMINAN EN UNA
% DIRECCIÓN («Disponible en: https://…» SIN PUNTO)
% ------------------------------------------------------------
\renewcommand*{\finentrypunct}{%
  \iffieldundef{doi}{\iffieldundef{url}{\addperiod}{}}{}}

% ------------------------------------------------------------
% FECHA: SOLO EL AÑO, COMO EN EL PATRÓN DE LA GUÍA. vancouver.bbx
% ARMA EL MES CON ABREVIATURAS EN INGLÉS (\shortmonth: Jan, Feb…)
% QUE EN UN ARTÍCULO EN CASTELLANO SALDRÍAN MEZCLADAS
% ------------------------------------------------------------
\DeclareFieldFormat*{date}{\thefield{year}}

% ------------------------------------------------------------
% DIRECCIONES: vancouver.bbx TRAE «Available from:» Y
% «[Accessed on: …]» FIJOS EN INGLÉS. LA GUÍA: «Disponible en:» Y
% «[consultado el 11 de junio de 2020]», CON LA FECHA LARGA DEL
% IDIOMA (LA QUE ARMA biblatex, SIN ABREVIAR: dateabbrev=false;
% EL AÑO DE LA REFERENCIA NO SE ABREVIA NUNCA)
% ------------------------------------------------------------
\DeclareFieldFormat*{url}{Disponible en:\space\url{#1}}
\DeclareFieldFormat*{urldate}{\mkbibbrackets{consultado el\space#1}}

% ------------------------------------------------------------
% LIBRO: «Lugar: Editorial; año» (LA GUÍA: «Madrid: Elsevier
% Science; 2007»). vancouver.bbx DEJA LA COMA DE biblatex
% ------------------------------------------------------------
\renewbibmacro*{publisher+location+date}{%
  \printlist{location}%
  \iflistundef{publisher}
    {\setunit*{\addcomma\space}}
    {\setunit*{\addcolon\space}}%
  \printlist{publisher}%
  \setunit*{\addsemicolon\space}%
  \usebibmacro{date}%
  \newunit}

% ------------------------------------------------------------
% CAPÍTULO: «Autor. Capítulo. En: Editor A, editor. Libro. Lugar:
% Editorial; año. p. 465-78.» (LA GUÍA). LOS EDITORES VAN DESPUÉS DE
% «En:» Y NO DESPUÉS DEL LIBRO; LAS PÁGINAS CON «p.», QUE EN
% VANCOUVER NO SE TRADUCE
% ------------------------------------------------------------
\renewbibmacro*{in:}{%
  \ifboolexpr{test {\ifentrytype{incollection}} or test {\ifentrytype{inbook}}
               or test {\ifentrytype{inproceedings}}}
    {\bibstring{in}\addcolon\space
     \ifnameundef{editor}{}
       {\printnames{editor}\addcomma\space
        \ifnumgreater{\value{editor}}{1}{\bibstring{gbeditores}}{\bibstring{gbeditor}}%
        \setunit{\addperiod\space}}}
    {}}
% «editor»/«editores» COMO LA GUÍA: spanish.lbx DA «edición» A
% PROPÓSITO (EVITA EL GÉNERO; LEÍDO EN spanish.lbx), Y ESE ES EL
% CRITERIO DE LOS OTROS ESTILOS, QUE NO SE TOCA. ESTAS DOS CADENAS SON
% SOLO DE VANCOUVER
\NewBibliographyString{gbeditor,gbeditores}
\DefineBibliographyStrings{spanish}{gbeditor = {editor}, gbeditores = {editores}}
\DefineBibliographyStrings{english}{gbeditor = {editor}, gbeditores = {editors}}
\letbibmacro{gb:byeditor+others}{byeditor+others}
\renewbibmacro*{byeditor+others}{%
  \ifboolexpr{test {\ifentrytype{incollection}} or test {\ifentrytype{inbook}}
               or test {\ifentrytype{inproceedings}}}
    {}
    {\usebibmacro{gb:byeditor+others}}}
\DeclareFieldFormat*[incollection,inbook,inproceedings]{pages}{%
  \adddot\addspace p.\addnbspace\mkcomprange{#1}}

